\documentclass[11pt, a4paper]{article}
\usepackage[margin=2.5cm]{geometry}
\usepackage{kotex}
\usepackage{amsmath, amssymb, mathtools}
\usepackage{amsthm}
\usepackage{booktabs}
\usepackage{xcolor}
\usepackage{hyperref}
\usepackage{listings}
\usepackage{tikz}
\usetikzlibrary{arrows.meta, positioning, shapes.geometric, calc}
\usepackage{tcolorbox}
\tcbuselibrary{skins, breakable}

\definecolor{qcDarkBlue}{HTML}{1A365D}
\definecolor{qcAccentBlue}{HTML}{2B6CB0}
\definecolor{qcLightBlue}{HTML}{EBF8FF}

\hypersetup{
    colorlinks=true,
    linkcolor=qcDarkBlue,
    urlcolor=qcAccentBlue,
    citecolor=qcDarkBlue
}

\newcommand{\ket}[1]{|#1\rangle}
\newcommand{\bra}[1]{\langle#1|}
\newcommand{\braket}[2]{\langle#1|#2\rangle}
\newcommand{\ketbra}[2]{|#1\rangle\langle#2|}

\lstset{
    basicstyle=\small\ttfamily,
    keywordstyle=\color{qcAccentBlue}\bfseries,
    commentstyle=\color{gray},
    stringstyle=\color{teal},
    backgroundcolor=\color{qcLightBlue!30},
    frame=single,
    rulecolor=\color{gray!30},
    breaklines=true,
    tabsize=4,
    showstringspaces=false
}

\theoremstyle{definition}
\newtheorem{definition}{정의}[section]
\newtheorem{theorem}{정리}[section]
\newtheorem{example}{예제}[section]

\newtcolorbox{notebox}[1]{
    colback=qcLightBlue!50,
    colframe=qcDarkBlue,
    fonttitle=\bfseries\color{white},
    title=#1,
    arc=2mm,
    left=8pt, right=8pt, top=8pt, bottom=8pt
}

\newtcolorbox{cautionbox}[1]{
    colback=red!5,
    colframe=red!75!black,
    fonttitle=\bfseries\color{white},
    title=#1,
    arc=2mm,
    left=8pt, right=8pt, top=8pt, bottom=8pt
}

\newtcolorbox{worksheetbox}[1]{
    colback=blue!3,
    colframe=qcAccentBlue,
    fonttitle=\bfseries\color{white},
    title=#1,
    arc=2mm,
    left=8pt, right=8pt, top=8pt, bottom=8pt
}

\title{\textbf{\Huge ‘양자컴퓨팅 이론\&실습 4주 부트캠프’}\\[0.4em]
\LARGE Day 2 공식 강의노트: 다중 큐비트, 텐서곱, 그리고 양자 얽힘 (NumPy 실습)\\[0.2em]
\large 선형대수학으로 정립하는 다차원 복소 힐베르트 공간과 비분리 얽힘 상태}
\author{\textbf{강사}: 서상협 ((주)큐렌스 / QRens)}
\date{2026년 11월 18일}

\begin{document}

\maketitle

\begin{abstract}
본 강의노트는 4주 부트캠프의 2일차(Day 2) 공식 교재입니다. 단일 큐비트의 2차원 상태 공간을 넘어 텐서곱(Tensor Product / Kronecker Product)을 통해 $2^n$ 차원의 복소 힐베르트 공간($\mathbb{C}^4$)으로 확장하는 원리를 정립합니다. 고전 컴퓨터의 독립된 비트 결합과 구별되는 큐비트 얽힘(Quantum Entanglement)의 본질, 2큐비트 상호작용 CNOT 게이트의 $4 \times 4$ 행렬 연산자, 최대 얽힘 상태인 4대 Bell 기저의 대수적 합성, 비분리성(Non-separability) 판별식($c_{00}c_{11}-c_{01}c_{10} \neq 0$), 그리고 부분 측정(Partial Measurement) 시 즉각적 파동함수 사영 붕괴와 No-Signaling 정리를 체계적으로 다룹니다.
\end{abstract}

\vspace{0.3cm}

\begin{notebox}{Day 2 학습의 3대 핵심 기둥}
\begin{enumerate}
    \item \textbf{[차원의 지수적 폭발과 텐서곱]}: $n$개 큐비트가 결합할 때 상태 공간의 차원은 덧셈($2n$)이 아닌 곱셈($2^n$)으로 지수 팽창하며, 다중 큐비트의 상태와 연산자는 Kronecker 텐서곱($A \otimes B$) 및 혼합곱 공식 $(A \otimes B)(C \otimes D) = (AC) \otimes (BD)$으로 정확히 기술됩니다.
    \item \textbf{[양자 우위의 핵심 자원, 비분리 얽힘]}: 얽힘 상태는 개별 단일 큐비트의 곱상태로 환원될 수 없는($c_{00}c_{11} - c_{01}c_{10} \neq 0$) 전역적 상관관계이며, 고전 컴퓨터로는 시뮬레이션하기 어려운 양자 알고리즘 가속의 원천입니다.
    \item \textbf{[부분 측정과 비국소 상관성]}: 얽힌 2큐비트 쌍에서 한 큐비트를 관측하는 순간 다른 큐비트의 상태가 즉각적으로 사영 붕괴하지만, 이 상관관계는 국소적 무작위성으로 인해 정보를 초광속으로 전송할 수 없습니다(No-Signaling 정리).
\end{enumerate}
\end{notebox}

% =============================================================================
\section{텐서곱(Tensor Product)과 \texorpdfstring{$2^n$}{2^n} 차원 복소 힐베르트 공간}
% =============================================================================
두 개의 단일 큐비트 계 $\mathcal{H}_1 \cong \mathbb{C}^2$과 $\mathcal{H}_2 \cong \mathbb{C}^2$가 결합된 복합 양자계는 텐서곱 공간으로 정의됩니다:
$$\mathcal{H} = \mathcal{H}_1 \otimes \mathcal{H}_2 \cong \mathbb{C}^4$$
$n$개의 큐비트가 결합하면 상태 공간의 차원은 $2^n$으로 지수 폭발합니다.

고전 컴퓨터의 50비트는 메모리에 0과 1 상태 50개만 기록하면 되지만, 50큐비트의 일반적인 중첩 상태벡터는 $2^{50} \approx 1.125 \times 10^{15}$개의 복소 진폭을 필요로 합니다.
각 복소수를 배정밀도(16바이트)로 저장할 경우 약 16페타바이트(PB)의 RAM이 필요하며, 이는 현존 최고 성능의 슈퍼컴퓨터 메모리를 즉시 초과합니다. 큐비트가 증가함에 따라 고전 컴퓨터가 겪는 이 '차원의 저주(Curse of Dimensionality)'가 역설적으로 양자컴퓨터가 고전 한계를 돌파하는 '양자 우위(Quantum Advantage)'의 배경이 됩니다.

\begin{definition}[텐서곱과 혼합곱 성질]
두 열벡터의 텐서곱 (Kronecker Product):
$$\ket{a} \otimes \ket{b} = \begin{pmatrix} a_0 \\ a_1 \end{pmatrix} \otimes \begin{pmatrix} b_0 \\ b_1 \end{pmatrix} = \begin{pmatrix} a_0 \begin{pmatrix} b_0 \\ b_1 \end{pmatrix} \\ a_1 \begin{pmatrix} b_0 \\ b_1 \end{pmatrix} \end{pmatrix} = \begin{pmatrix} a_0 b_0 \\ a_0 b_1 \\ a_1 b_0 \\ a_1 b_1 \end{pmatrix}$$
표준 4차원 계산 기저벡터는 다음과 같이 구성됩니다:
$$\ket{00} = \begin{pmatrix} 1 \\ 0 \\ 0 \\ 0 \end{pmatrix}, \quad \ket{01} = \begin{pmatrix} 0 \\ 1 \\ 0 \\ 0 \end{pmatrix}, \quad \ket{10} = \begin{pmatrix} 0 \\ 0 \\ 1 \\ 0 \end{pmatrix}, \quad \ket{11} = \begin{pmatrix} 0 \\ 0 \\ 0 \\ 1 \end{pmatrix}$$

\textbf{행렬과 벡터의 혼합곱 공식(Mixed-Product Property)}:
$$(A \otimes B)(C \otimes D) = (AC) \otimes (BD)$$
이 성질은 개별 큐비트에 독립적으로 게이트를 가하는 연산이 텐서곱 공간 전체에서 어떻게 교환·합성되는지를 대수적으로 규정합니다.
\end{definition}

% =============================================================================
\section{곱상태(Product State)와 얽힘 상태(Entangled State)}
% =============================================================================
일반적인 2큐비트 순수 상태는 4차원 복소 단위벡터로 표현됩니다:
$$\ket{\psi} = c_{00}\ket{00} + c_{01}\ket{01} + c_{10}\ket{10} + c_{11}\ket{11}, \quad \sum_{x \in \{00,01,10,11\}} |c_x|^2 = 1$$

\begin{definition}[곱상태와 비분리성 판별식]
\begin{itemize}
    \item \textbf{곱상태(Product State / Separable State)}:
    상태벡터가 두 단일 큐비트 상태의 단순 텐서곱으로 인수분해될 수 있는 상태:
    $$\ket{\psi} = (\alpha_0\ket{0} + \alpha_1\ket{1}) \otimes (\beta_0\ket{0} + \beta_1\ket{1}) = \alpha_0\beta_0\ket{00} + \alpha_0\beta_1\ket{01} + \alpha_1\beta_0\ket{10} + \alpha_1\beta_1\ket{11}$$
    이때 계수 행렬의 행렬식(Determinant)은 항상 0입니다:
    $$\det \begin{pmatrix} c_{00} & c_{01} \\ c_{10} & c_{11} \end{pmatrix} = c_{00}c_{11} - c_{01}c_{10} = (\alpha_0\beta_0)(\alpha_1\beta_1) - (\alpha_0\beta_1)(\alpha_1\beta_0) = 0$$
    \item \textbf{얽힘 상태(Entangled State)}:
    두 큐비트가 상호작용하여 상태벡터가 어떠한 단일 큐비트들의 텐서곱으로도 인수분해될 수 없는 비분리(Non-separable) 상태:
    $$c_{00}c_{11} - c_{01}c_{10} \neq 0$$
\end{itemize}
\end{definition}

컴퓨터과학적 관점에서 얽힘이 없는 상태는 50큐비트라 할지라도 100개의 복소 진폭만으로 일반 스마트폰에서 수 밀리초 만에 완벽히 시뮬레이션할 수 있습니다. 즉, 얽힘이 결여된 양자 회로는 고전 컴퓨터 대비 계산 우위를 제공하지 못하며, \textbf{비분리 얽힘이야말로 양자 우위의 유일한 원천}입니다.

% =============================================================================
\section{CNOT 게이트와 Bell 상태의 대수적 합성}
% =============================================================================
다중 큐비트 사이의 얽힘을 생성하기 위해서는 큐비트 간 상호작용을 매개하는 2큐비트 유니터리 연산자가 필요합니다.
대표적인 연산자가 \textbf{CNOT(Controlled-NOT, $CX$)} 게이트입니다.

CNOT 게이트는 제어 큐비트($q_0$)가 $\ket{1}$일 때 타겟 큐비트($q_1$)를 반전($X$)시키고, 제어 큐비트가 $\ket{0}$일 때는 타겟 큐비트를 변경하지 않는 $4 \times 4$ 유니터리 연산자입니다:
$$CX = \begin{pmatrix} 1 & 0 & 0 & 0 \\ 0 & 1 & 0 & 0 \\ 0 & 0 & 0 & 1 \\ 0 & 0 & 1 & 0 \end{pmatrix}$$

\begin{worksheetbox}{[손계산 워크시트] 2큐비트 Bell 상태 ($\ket{\Phi^+}$) 단계별 수식 유도}
\textbf{문제}: 초기 상태 $\ket{\psi_0} = \ket{00}$에 Hadamard 게이트와 CNOT 게이트를 순차적으로 적용할 때, 상태벡터의 변화를 $4 \times 4$ 행렬과 $4 \times 1$ 열벡터 연산으로 직접 유도하시오.

\textbf{[단계별 유도 과정]}:
\begin{enumerate}
    \item \textbf{초기 상태의 4차원 벡터 표현}:
    $$\ket{\psi_0} = \ket{0} \otimes \ket{0} = \begin{pmatrix} 1 \\ 0 \end{pmatrix} \otimes \begin{pmatrix} 1 \\ 0 \end{pmatrix} = \begin{pmatrix} 1 \\ 0 \\ 0 \\ 0 \end{pmatrix} = \ket{00}$$
    \item \textbf{1단계: $q_0$에 $H$ 게이트 작용 ($H \otimes I$ 연산자)}:
    $$\ket{\psi_1} = (H \otimes I)\ket{00} = (H\ket{0}) \otimes \ket{0} = \left(\frac{\ket{0}+\ket{1}}{\sqrt{2}}\right) \otimes \ket{0} = \frac{\ket{00} + \ket{10}}{\sqrt{2}} = \frac{1}{\sqrt{2}}\begin{pmatrix} 1 \\ 0 \\ 1 \\ 0 \end{pmatrix}$$
    이 상태는 $c_{00}c_{11} - c_{01}c_{10} = (1/\sqrt{2})(0) - (0)(1/\sqrt{2}) = 0$이므로 여전히 \textbf{분리 가능한 곱상태}입니다.
    \item \textbf{2단계: CNOT 게이트 작용 (제어: $q_0$, 타깃: $q_1$)}:
    $$\ket{\psi_2} = CX \ket{\psi_1} = \begin{pmatrix} 1 & 0 & 0 & 0 \\ 0 & 1 & 0 & 0 \\ 0 & 0 & 0 & 1 \\ 0 & 0 & 1 & 0 \end{pmatrix} \begin{pmatrix} 1/\sqrt{2} \\ 0 \\ 1/\sqrt{2} \\ 0 \end{pmatrix} = \begin{pmatrix} 1/\sqrt{2} \\ 0 \\ 0 \\ 1/\sqrt{2} \end{pmatrix} = \frac{\ket{00} + \ket{11}}{\sqrt{2}} = \mathbf{\ket{\Phi^+}}$$
    \item \textbf{얽힘 판별식 검증}:
    $$c_{00}c_{11} - c_{01}c_{10} = \left(\frac{1}{\sqrt{2}}\right)\left(\frac{1}{\sqrt{2}}\right) - (0)(0) = \frac{1}{2} \neq 0$$
    따라서 완벽한 비분리 \textbf{최대 얽힘 상태(Maximally Entangled State)}가 완성되었습니다!
    \item \textbf{측정 확률}:
    $$P(00) = \left|\frac{1}{\sqrt{2}}\right|^2 = 50\%, \qquad P(11) = \left|\frac{1}{\sqrt{2}}\right|^2 = 50\%, \qquad P(01) = P(10) = 0\%$$
\end{enumerate}
\end{worksheetbox}

\subsection{4대 Bell 기저 (Bell Basis)}
2큐비트 힐베르트 공간 $\mathbb{C}^4$의 기저를 이루는 4가지 정규직교 최대 얽힘 상태는 다음과 같이 정의됩니다:
\begin{align*}
\ket{\Phi^+} &= \frac{\ket{00} + \ket{11}}{\sqrt{2}}, \qquad \ket{\Phi^-} = \frac{\ket{00} - \ket{11}}{\sqrt{2}} \\[0.5em]
\ket{\Psi^+} &= \frac{\ket{01} + \ket{10}}{\sqrt{2}}, \qquad \ket{\Psi^-} = \frac{\ket{01} - \ket{10}}{\sqrt{2}}
\end{align*}
이 4가지 상태는 상호 직교($\braket{\Phi^\pm}{\Psi^\pm} = 0$)하며, 양자 원격전송(Teleportation), 초고밀도 부호화(Superdense Coding), 양자 암호통신(E91 프로토콜)의 근간이 됩니다.

% =============================================================================
\section{부분 측정(Partial Measurement)과 No-Signaling 정리}
% =============================================================================
Bell 상태 $\ket{\Phi^+} = \frac{\ket{00}+\ket{11}}{\sqrt{2}}$에서 첫 번째 큐비트($q_0$)만을 측정하는 경우를 분석합니다.

\begin{enumerate}
    \item \textbf{$q_0$ 측정 결과의 확률}:
    $$P(q_0 = 0) = |c_{00}|^2 + |c_{01}|^2 = \left|\frac{1}{\sqrt{2}}\right|^2 + 0 = 50\%$$
    $$P(q_0 = 1) = |c_{10}|^2 + |c_{11}|^2 = 0 + \left|\frac{1}{\sqrt{2}}\right|^2 = 50\%$$
    \item \textbf{사영 붕괴 후 $q_1$의 상태}:
    \begin{itemize}
        \item $q_0 = 0$이 관측된 경우: 복합계 상태는 즉각 $\ket{00}$으로 투영 붕괴하므로, 두 번째 큐비트는 아무런 조작을 가하지 않았음에도 \textbf{$100\%$ 확률로 $\ket{0}$}으로 결정됩니다.
        \item $q_0 = 1$이 관측된 경우: 복합계 상태는 즉각 $\ket{11}$으로 붕괴하므로, 두 번째 큐비트는 \textbf{$100\%$ 확률로 $\ket{1}$}로 결정됩니다.
    \end{itemize}
\end{enumerate}

\begin{cautionbox}{물리적 성찰: 아인슈타인의 EPR 역설과 No-Signaling 정리}
두 큐비트가 은하계 양 끝으로 수십억 광년 떨어져 있더라도, 한쪽에서 $q_0$를 측정하면 반대편 $q_1$의 상태는 그 즉시 결정됩니다. 아인슈타인은 이를 "유령 같은 원격 작용(Spooky action at a distance)"이라 부르며 양자역학의 불완전성을 주장했습니다.

그러나 $q_1$을 가진 관측자 입장에서는 $q_0$의 측정 여부를 미리 알 수 없으며, $q_1$ 단독 측정 결과는 언제나 무작위 $50\%:50\%$로만 나타납니다. 두 관측자가 결과를 대조(상관관계 확인)하기 위해서는 반드시 빛의 속도 이하로 고전 통신을 주고받아야 하므로, \textbf{초광속 정보 전달(Superluminal Signaling)은 엄격히 불가능}합니다(No-Signaling 정리).
\end{cautionbox}

\end{document}
